\documentclass[journal=esthag,manuscript=article,margin=1in]{achemso}
\usepackage[T1]{fontenc}
\usepackage{amsmath,amssymb}
\usepackage{graphicx}
\usepackage{float}
\usepackage{url}
\setlength{\emergencystretch}{2em}

\author{Jeongyun Kim}
\altaffiliation{These authors contributed equally to this work.}
\affiliation[SNU]{School of Chemical and Biological Engineering, Seoul National University, Seoul, Republic of Korea}
\author{Jin Yoo}
\altaffiliation{These authors contributed equally to this work.}
\affiliation[CBDRI]{CBRN Defense Research Institute, Republic of Korea}
\author{Ku Kang}
\email{bisu9082@gmail.com}
\affiliation[CBDRI]{CBRN Defense Research Institute, Republic of Korea}

\title[Conditional uncertainty for PFAS clearance]
{Can Structure-Based Models Clear Untested PFAS? External Validation Reveals Endpoint-, Representation-, and Distribution-Dependent Reliability}

\keywords{PFAS; aquatic ecotoxicity; conformal prediction; distribution shift; uncertainty quantification; QSAR; applicability domain}

\begin{document}

\begin{tocentry}
\includegraphics[width=3.25in,height=1.75in,keepaspectratio]{fig_02_framework.png}
\end{tocentry}

\section*{Synopsis}
Machine-learning intervals calibrated over general chemicals can be misleading for PFAS; PFAS-conditional calibration restores held-out coverage but remains too wide to support low-hazard clearance under the tested protective endpoint.

\begin{abstract}
Structure-based screening could help prioritize the large PFAS inventory, but a model useful for prioritization is not automatically reliable for clearing untested substances as low concern. We evaluated that distinction using freshwater aquatic-toxicity records from EnviroTox and an external set of 138 measured PFAS. A structure-locked analysis removed PFAS matches and duplicate desalted structures, leaving 3,788 unique non-PFAS training structures. Random-forest models using ECFP4 fingerprints performed moderately in repeated broad-space cross-validation ($R^2=0.39\pm0.02$) but failed on fully unseen PFAS ($R^2=-0.58$; Spearman $\rho=-0.07$). At nominal 90\% coverage, intervals calibrated on broad non-PFAS chemicals covered only 0.659 of PFAS on average, whereas PFAS-conditional split calibration covered 0.922. The gain in held-out coverage required much wider intervals (mean median width 10.53 $\log_{10}$ mg/L). After removing 22 unparsable records, 140 records overlapping measured PFAS, and duplicate structures from the 19,028-record inventory, 18,851 unique untested structures remained; 17,128 passed the prespecified applicability-domain filter. Broad marginal intervals would clear 3,138 in-domain structures, despite their demonstrated PFAS undercoverage. PFAS-conditional intervals cleared none. This zero is conditional on the measured-to-untested PFAS exchangeability assumption and does not imply that every PFAS is hazardous. Endpoint and representation mattered: on 97 PFAS with acute data, acute-only rank prediction improved relative to the most-sensitive mixed target (paired $\Delta\rho=+0.25$, 95\% bootstrap CI $[+0.09,+0.42]$), and 165 physicochemical descriptors restored positive rank prediction for the protective target ($\rho=+0.42$); in every case absolute accuracy stayed low ($R^2$ near zero), an improvement in ranking rather than clearance-grade prediction. The results support prioritization, not safety certification: predictive validity must be demonstrated for the target chemical population, endpoint, and representation before model scores are used for clearance.
\end{abstract}

\section{Introduction}

Per- and polyfluoroalkyl substances (PFAS) encompass a structurally and functionally diverse chemical universe whose boundaries depend on the operational definition used \cite{buck2011terminology,oecd2021pfas}. Their persistence, widespread occurrence, and regulatory relevance have motivated broad inventory and risk-management programs \cite{wang2017never,cousins2022planetary,epa2021roadmap}. Human-exposure and environmental-risk reviews further show that PFAS cannot be treated as a toxicologically uniform class \cite{sunderland2019exposure,brendel2018short,gomis2018potency,conder2008pfca}. Chain length, functional group, and protein or receptor interactions can alter biological behavior \cite{cheng2018protein,wolf2008ppar,vandenheuvel2006ppar}. Consequently, the practical problem is not simply to assign a score to every listed structure, but to determine which decisions those scores can support.

The U.S. EPA CompTox Chemicals Dashboard and its PFASSTRUCTV6 list provide a large structure-based inventory \cite{williams2017comptox,epacomptox2026}. By contrast, measured aquatic effects cover only a small subset. EnviroTox and ECOTOX consolidate valuable ecotoxicity evidence, but assay coverage remains heterogeneous across species, duration, and effect type \cite{connors2019envirotox,olker2022ecotox}. Category-based testing strategies therefore remain necessary \cite{patlewicz2019category}. Structure standardization is equally important because salts, isotopologues, and alternative registry forms can otherwise inflate apparent sample size or create leakage between model partitions \cite{mceachran2018msready}.

Quantitative structure--activity relationship models are natural tools for filling these gaps. Aquatic-toxicity QSARs have a long history \cite{mazzatorta2005aquatic,nath2023aquatic}, and current environmental applications combine curated descriptors, fingerprints, and machine learning \cite{mansouri2018opera,ms2tox2022}. Best-practice guidance emphasizes chemically independent validation, transparent applicability domains, and an explicit definition of the modeled endpoint \cite{tropsha2010best,cherkasov2014qsar}. Our prior work used machine learning for multi-endpoint environmental prioritization and mechanistic toxicity modeling \cite{kim2026dualthreat,jang2025opqsar}; the present study asks the narrower and more demanding question of whether such models can support PFAS \emph{clearance}. This is more than a modeling question. Structure-based scores are already generated for entire PFAS inventories, and an unexamined step from ``low predicted hazard'' to ``de-prioritized for testing'' would convert a prioritization aid into a \emph{de facto} clearance decision. Establishing when that step is, and is not, statistically defensible is therefore a safeguard against regulatory misuse of machine-learning screening, not merely a model-validation exercise.

Conformal prediction can attach finite-sample prediction intervals under exchangeability \cite{vovk2005algorithmic,shafer2008tutorial,angelopoulos2023gentle}. It has been proposed as a transparent uncertainty framework for chemical prediction \cite{norinder2014conformal}. The guarantee, however, is marginal over the calibration distribution and need not hold for a shifted subgroup. Mondrian or group-conditional calibration can target a prespecified subgroup, but its interpretation still depends on exchangeability within that group \cite{vovk2012conditional}. This distinction is central when measured PFAS are selectively tested and may not represent the untested inventory.

Here we separate four questions: (1) does a broad aquatic-toxicity model perform on general chemicals; (2) does it transfer to fully unseen PFAS; (3) do its intervals cover held-out measured PFAS; and (4) are those intervals informative enough to clear untested structures above a protective threshold? We also test whether endpoint construction and molecular representation govern transfer. The analysis is deliberately conservative: duplicate structures and measured-inventory overlap are removed before model fitting and deployment, and uncertainty claims are limited to the validation design actually tested.

\section{Materials and Methods}

\subsection{Data sources, endpoint, and structure lock}

EnviroTox v2.0 contains 80,912 effect records and 4,267 substance entries \cite{connors2019envirotox}. We retained freshwater fish, invertebrate, algae, and amphibian records reported in mass concentration units. Replicate values were combined by the geometric mean within trophic-level and endpoint-class strata. The primary protective target was the minimum concentration across acute LC50/EC50 and chronic NOEC/LOEC/MATC strata, modeled as $y=\log_{10}(T/\mathrm{mg\,L^{-1}})$.

RDKit \cite{rdkit} parsed 4,068 broad chemical records. PFAS were removed by CAS or canonical-structure matching to PFASSTRUCTV6 and to the measured-PFAS set. Because 213 rows duplicated a desalted canonical structure, identical structures were locked to one modeling row using the minimum $y$, consistent with the protective target. The final broad training set comprised 3,788 unique non-PFAS structures, including 936 at or below 0.1 mg/L. The independent PFAS set comprised 138 unique measured structures, of which 57 were at or below 0.1 mg/L.

PFASSTRUCTV6 contained 19,028 inventory records in the supplied snapshot. Before deployment, we removed 22 records whose SMILES could not be parsed, 140 records that overlapped measured PFAS by CAS or canonical structure, and duplicate canonical structures. This produced 18,851 unique, parsable, operationally untested structures. This count is specific to the PFASSTRUCTV6 operational definition and preprocessing rule; it is not asserted as a universal count of PFAS.

\subsection{Representations and predictive models}

The default representation was a 2,048-bit ECFP4 fingerprint (Morgan radius 2) \cite{rogers2010ecfp}. We also evaluated hashed Morgan counts and 165 finite RDKit physicochemical descriptors after excluding Ipc and median-imputing missing values. The default learner was a random-forest regressor with 300 trees, $\texttt{max\_features}=\sqrt{p}$, and random seed 20260620 \cite{breiman2001random,sklearn}. Broad-space performance used five repeats of shuffled five-fold cross-validation over the structure-locked rows. The external PFAS model was fitted to all 3,788 non-PFAS structures and evaluated once on the 138 PFAS.

For endpoint sensitivity, acute-only and chronic-only targets were reconstructed from the same source records. Separate non-PFAS models were fitted for each endpoint. The acute comparison included 97 measured PFAS and the chronic comparison 108. The paired endpoint effect compared acute-only and most-sensitive predictions on the same 97 PFAS. Spearman correlations were bootstrapped 5,000 times at the substance level.

\subsection{Conformal validation}

Normalized nonconformity was
\begin{equation}
s_i=\frac{|y_i-\hat y_i|}{\sigma_i+\beta},
\end{equation}
where $\sigma_i$ is the standard deviation across forest-tree predictions and $\beta$ is the median calibration-set $\sigma$ plus $10^{-6}$. Quantiles used the finite-sample higher-order rule $\lceil(n+1)(1-\alpha)\rceil/n$.

Two calibration designs were evaluated over 40 prespecified random repeats. For marginal calibration, 20\% of broad non-PFAS structures calibrated a model fitted to the remaining 80\%, and the resulting intervals were evaluated on all 138 PFAS. For PFAS-conditional calibration, the model remained fitted only to broad non-PFAS structures while the 138 PFAS were split into disjoint calibration and test halves. We report the mean coverage and the 2.5th--97.5th percentiles across repeats rather than treating overlapping repeated-split predictions as independent Bernoulli trials.

For inventory deployment, the final broad model was calibrated using all 138 measured PFAS. This produces a PFAS-conditional interval under the explicit assumption that measured and untested PFAS are exchangeable after conditioning on PFAS membership. We do not claim that the assumption is guaranteed. A structure was \emph{clearable} when the lower (more hazardous) bound of its 90\% interval exceeded 0.1 mg/L.

\subsection{Applicability domain and software}

The applicability domain was the mean Euclidean distance in binary ECFP4 space to the five nearest non-PFAS training structures. The threshold was the 95th percentile of leave-self-out training distances. Analyses used Python, scikit-learn, NumPy, SciPy, and RDKit \cite{sklearn,rdkit}. Processed, structure-locked tables, split-level results, and figure-generating scripts accompany the revised package. Future sequential testing should benchmark simple hazard-ranking baselines against pool-based active-learning approaches rather than assuming that uncertainty sampling is superior \cite{settles2009active,graff2021molpal,smith2018less,lookman2019active}.

\section{Results and Discussion}

\subsection{Data locking changes the denominator, not the scientific question}

The original chemical-level table contained 4,068 parsable records but only 3,855 unique canonical structures before PFAS removal. Seventy-three broad rows matched the PFAS inventory or measured PFAS, leaving 3,788 unique non-PFAS training structures after structure locking (Figure~\ref{fig:data}). The inventory audit was also consequential: the supplied 19,028 records were not all untested unique structures. Removing invalid structures, measured overlap, and structural duplicates yielded 18,851 deployment structures. These corrections prevent the same molecule from being counted repeatedly or appearing on both sides of the measured--untested boundary.

\begin{figure*}[t]
\centering
\includegraphics[width=\linewidth]{fig_01_data_audit.png}
\caption{Auditable data lock. (a) Most-sensitive toxicity distribution for 3,788 unique non-PFAS training structures. (b) Repeated broad-space cross-validation. (c) Inventory audit separating invalid SMILES and measured-PFAS overlap from the final unique untested set. (d) Record-to-structure flow used to prevent structural leakage.}
\label{fig:data}
\end{figure*}

\begin{figure*}[t]
\centering
\includegraphics[width=\linewidth]{fig_02_framework.png}
\caption{Validation framework. Model training, external PFAS validation, group-conditional calibration, and inventory deployment are treated as distinct stages. The final deployment statement is explicitly conditional on exchangeability between measured and untested PFAS.}
\label{fig:framework}
\end{figure*}

\subsection{Broad-space accuracy does not transfer to PFAS}

The structure-locked ECFP4 model achieved mean fold $R^2=0.39\pm0.02$ in 25 broad-space cross-validation folds. This result establishes moderate interpolation performance for general non-PFAS aquatic-toxicity data; it does not establish PFAS transfer. On the 138 fully unseen PFAS, $R^2$ fell to $-0.58$ and Spearman $\rho$ to $-0.07$ (Figure~\ref{fig:validity}a). Predictions contracted toward the center of the training response distribution and failed to rank PFAS toxicity. For a screening program this is the decisive gap: a model that looks trustworthy in cross-validation can fail precisely on the chemical class a clearance decision would target. The failure was not specific to random forests---gradient-boosting and ridge regression on the same fingerprints also failed externally (Spearman $\rho=-0.04$ and $+0.09$, both $R^2<0$; Supporting Information)---so it reflects the representation and the distribution shift rather than one algorithm.

The contrast also clarifies why a fold-level increase in broad-space $R^2$ must not be described as a fold improvement in PFAS accuracy. The two evaluations answer different questions. The revised analysis therefore reports broad interpolation and PFAS external validation separately rather than presenting their ratio as an accuracy gain.

\subsection{Marginal intervals under-cover PFAS}

At nominal 90\% coverage, broad marginal calibration covered PFAS with mean probability 0.659 across the 40 repeats (2.5th--97.5th percentile 0.623--0.703). PFAS-conditional split calibration increased mean held-out coverage to 0.922 (0.841--0.986; Figure~\ref{fig:validity}b). This is evidence that group-aware calibration can restore held-out coverage for measured PFAS under the tested resampling design.

The improvement was not free. At the same nominal level, the mean median interval width increased from 4.84 to 10.53 $\log_{10}$ mg/L (Figure~\ref{fig:validity}c). No known-low-hazard PFAS in the held-out halves was cleared at nominal coverage levels of 0.80 or higher. Thus coverage and decision power must be reported together: an interval can be empirically valid and still be too wide to support a useful clearance decision.

\begin{figure}[H]
\centering
\includegraphics[width=\linewidth]{fig_03_validity.png}
\caption{External validity and conformal consequences. (a) Point prediction on 138 fully unseen PFAS. (b) Mean empirical coverage across 40 repeats; ribbons show the 2.5th--97.5th percentiles across repeats. (c) Corresponding median interval widths. (d) At 90\% nominal coverage, broad marginal intervals would clear 3,138 in-domain inventory structures despite demonstrated PFAS undercoverage, whereas PFAS-conditional intervals clear none under the measured-to-untested exchangeability assumption.}
\label{fig:validity}
\end{figure}

\subsection{No untested PFAS can be conditionally cleared at 90\%}

Of the 18,851 unique untested inventory structures, 17,128 (90.9\%) passed the prespecified fingerprint-distance domain. This membership was moderately threshold-dependent (15,037 at the 90th and 18,295 at the 99th training-distance percentile; Supporting Information), but the clearance conclusion below was not, because no in-domain structure cleared at any of these thresholds. A marginal cross-fitted calibration over broad chemicals produced median in-domain interval width 4.09 log units and would clear 3,138 structures. That number is not defensible as a PFAS safety decision because the same procedure covered only 0.659 of measured PFAS at the nominal 0.90 level.

Using all 138 measured PFAS for conditional deployment calibration increased the median in-domain width to 10.12 log units and cleared zero structures. The correct interpretation is limited but important: \emph{given the measured PFAS calibration sample, the selected protective endpoint, the ECFP4 model, and exchangeability between measured and untested PFAS, no inventory structure has enough evidence for 90\% low-hazard clearance}. This zero was stable under leave-one-PFAS-out recalibration (calibration-quantile coefficient of variation 0.8\%; Supporting Information), so it is not an artifact of a single calibration draw. The result neither proves that every PFAS is hazardous nor supplies an unconditional regulatory guarantee. Selective historical testing and structural heterogeneity remain plausible violations of measured-to-untested exchangeability.

\subsection{Endpoint and representation govern transfer}

The most-sensitive target mixes acute and chronic information according to assay availability. Under ECFP4, its external rank correlation was $-0.03$ in the per-configuration sensitivity fits (negligibly different from the single deployment fit above, $\rho=-0.07$; both indistinguishable from zero). Acute-only prediction was positive ($\rho=+0.30$, 95\% bootstrap CI $[+0.09,+0.48]$; $n=97$), whereas chronic-only prediction was near zero ($\rho=+0.02$; $n=108$; Figure~\ref{fig:endpoint}a). On the same 97 PFAS, changing from the mixed most-sensitive target to acute-only increased rank correlation by $+0.25$ (95\% CI $[+0.09,+0.42]$). The acute point $R^2$ remained slightly negative ($-0.06$), so this is evidence of improved ranking, not clearance-grade numeric accuracy.

Representation produced a second, independent sensitivity. Binary ECFP4 and Morgan counts both failed on the protective target ($\rho=-0.03$ and $-0.02$), whereas 165 RDKit physicochemical descriptors reached $\rho=+0.42$ (Figure~\ref{fig:endpoint}b). Descriptor-model $R^2$ remained $-0.06$, again indicating useful rank information without accurate absolute prediction. These results support the paper's central claim in a more precise form: PFAS transfer is a property of the complete endpoint--representation--calibration pipeline, not of an algorithm name alone.

\begin{figure}[H]
\centering
\includegraphics[width=\linewidth]{fig_04_endpoint_representation.png}
\caption{Endpoint and representation sensitivity after structure locking. (a) External Spearman correlation by endpoint; bars show point estimates and error bars show substance-level 95\% bootstrap intervals. (b) Representation sensitivity for the protective most-sensitive endpoint. (c) Acute-only external predictions. (d) Paired endpoint comparison on the same 97 PFAS.}
\label{fig:endpoint}
\end{figure}

\subsection{Limitations and implications}

First, the protective endpoint is assay-coverage dependent and combines biologically different effect measures. Its regulatory conservatism does not make it a homogeneous statistical target. Second, the measured PFAS set is modest and selectively observed. PFAS-conditional calibration addresses between-group shift only if the future PFAS are exchangeable with the calibration PFAS; membership in the same inventory does not guarantee this. Third, PFASSTRUCTV6 implements a broad operational definition and includes substances outside the narrower set often implied by ``legacy PFAS.'' Results should therefore be stratified by structural subgroup before laboratory prioritization. Fourth, per-tree dispersion is a pragmatic scale estimator rather than a mechanistic decomposition of uncertainty. We did not test learned molecular representations (e.g., graph neural networks); because the external failure is driven by distribution shift rather than model capacity---the training structures contain almost no fluorinated chemistry---a representation learned from the same data is expected to reproduce, not resolve, it. Finally, the study evaluates aquatic effect concentrations, not exposure, persistence, bioaccumulation, or overall risk.

These limitations do not erase the main result. They specify its proper use. The actionable, positive finding is specific: physicochemical descriptors on an endpoint-specific (acute) target were the only configuration to recover positive external rank prediction ($\rho=+0.42$--$0.48$). The model can therefore provide a usable provisional testing rank \emph{when those choices are made}, even though it cannot certify low hazard. It cannot presently certify absence of high aquatic hazard for an untested inventory structure at 90\% conditional confidence. Testing investments should target structural regions that are both poorly represented and decision-relevant, and each new measurement should be reserved appropriately between model updating and calibration.

\section{Conclusions}

A model that appears useful over general aquatic-toxicity data can fail both in point prediction and uncertainty coverage on PFAS. Structure locking and measured-inventory deduplication reduced the deployable inventory from 19,028 records to 18,851 unique untested structures. Marginal intervals calibrated on non-PFAS chemicals under-covered measured PFAS and would have cleared 3,138 in-domain structures. PFAS-conditional calibration restored held-out coverage but produced intervals too wide to clear any structure at 90\%. Endpoint and representation materially changed transfer, confirming that trustworthiness belongs to the complete modeling pipeline. Concretely, only physicochemical descriptors on an endpoint-specific target recovered positive rank prediction for the protective endpoint ($\rho=+0.42$), indicating a hypothesis-generating configuration for prioritization that warrants prospective validation before operational use. For present PFAS data, model scores are suitable for cautious prioritization, not low-hazard clearance.

\section*{Associated Content}
\textbf{Supporting Information.} Data-audit flow; repeated-split conformal results; endpoint and representation sensitivity; applicability-domain and algorithm robustness; and a numerical consistency manifest.

\section*{Data and Code Availability}
All processed and structure-locked input tables, split-level and canonical results, the numerical consistency manifest, and the analysis and figure-generation scripts that reproduce every reported number and figure are openly available at
\url{https://github.com/bisu9082/pfas-clearance-validation}.
The underlying source data are public: aquatic-toxicity records from the EnviroTox database and the U.S.\ EPA ECOTOX Knowledgebase, and the screening inventory from the U.S.\ EPA CompTox Chemicals Dashboard PFASSTRUCTV6 list. Raw third-party databases are not redistributed; the repository documents the exact versions and the reproducible processing pipeline used to regenerate the locked tables

\section*{Acknowledgement}
The authors thank the teams maintaining EnviroTox, U.S. EPA ECOTOX, and the CompTox Chemicals Dashboard.

\bibliography{myref}
\end{document}
